\documentclass[11pt]{report}

\usepackage[a4paper,margin=1in]{geometry}
\usepackage{fontspec}
\usepackage{titlesec}
\usepackage{enumitem}
\usepackage{longtable}
\usepackage{array}
\usepackage{booktabs}
\usepackage{xcolor}
\usepackage{hyperref}
\usepackage{tcolorbox}
\tcbuselibrary{skins,breakable}
\usepackage{amsmath,amssymb}

\setmainfont{TeX Gyre Pagella}

\definecolor{priorityA}{HTML}{1B4965}
\definecolor{priorityB}{HTML}{5F0F40}
\definecolor{priorityC}{HTML}{9A031E}
\definecolor{priorityD}{HTML}{0F4C5C}
\definecolor{priorityE}{HTML}{6A4C93}
\definecolor{priorityF}{HTML}{BC6C25}
\definecolor{priorityG}{HTML}{264653}
\definecolor{priorityH}{HTML}{606C38}

\titleformat{\chapter}[display]
  {\normalfont\huge\bfseries}{\chaptertitlename\ \thechapter}{20pt}{\Huge}
\titleformat{\section}{\normalfont\Large\bfseries}{\thesection}{1em}{}
\titleformat{\subsection}{\normalfont\large\bfseries}{\thesubsection}{1em}{}

\newtcolorbox{prioritybox}[2]{
  enhanced, breakable, colback=white, colframe=#1,
  boxrule=0.6pt, left=6pt, right=6pt, top=4pt, bottom=4pt,
  title={#2}, fonttitle=\bfseries\large, coltitle=white,
  colbacktitle=#1, attach boxed title to top left={yshift=-2mm,xshift=4mm},
  boxed title style={boxrule=0pt,colframe=#1}
}

\newtcolorbox{deliverable}{
  enhanced, breakable, colback=gray!5, colframe=gray!60,
  boxrule=0.4pt, left=6pt, right=6pt, top=3pt, bottom=3pt,
  fontupper=\small
}

\newcommand{\statuscheck}{\textcolor{green!50!black}{$\checkmark$}}
\newcommand{\statusgap}{\textcolor{red!70!black}{$\times$}}
\newcommand{\statuspartial}{\textcolor{orange!80!black}{$\sim$}}

% Macros matching the main manuscript's notation, used when
% referring to its theorems and definitions by their native symbols.
\newcommand{\reach}{\mathcal{R}}
\newcommand{\adm}{\mathcal{A}}
\newcommand{\repair}{\mathfrak{R}}
\newcommand{\reco}{\mathfrak{rec}}
\newcommand{\Vol}{\operatorname{Vol}}

\title{\Huge Editorial Completion Plan\\[0.3em]
  \Large \textit{The Ecology of Distinctions}\\[0.3em]
  \large From Theoretical Completeness to Pedagogical Completeness}
\author{}
\date{}

\begin{document}

\maketitle
\tableofcontents

\chapter*{Preface to the Plan}
\addcontentsline{toc}{chapter}{Preface to the Plan}

The manuscript, in its current state, comprises ten parts and
thirty-five chapters (thirty-one numbered chapters plus four
appendices), totaling approximately 13{,}250 lines of LuaLaTeX
source. The foundational chain

\begin{center}
\textit{Distinction} $\to$ \textit{Information} $\to$
\textit{Entropy} $\to$ \textit{History} $\to$
\textit{Recoverability} $\to$ \textit{Repair} $\to$
\textit{Regeneration} $\to$ \textit{Reachability} $\to$
\textit{Admissibility} $\to$ \textit{Generativity}
\end{center}

has been formalized end to end: every chapter contains
theorems with complete proofs, every cross-reference resolves,
and the dependency structure connecting Parts I--X is explicit
and machine-checked. In this sense the manuscript's
\emph{mathematical} architecture is complete.

What remains is a different kind of completeness. A reader
encountering Chapter 16 today is handed the RSVP field
equations within the first two pages, with no extended
discussion of why a scalar-vector-entropy triple is the
\emph{right} minimal field content, what alternatives were
rejected and why, or how a reader without a physics background
should approach the material. The same pattern recurs at every
``realization'' chapter: the abstract machinery of Parts I--III
is sound, but the bridge into each concrete domain is short.

This document is not a request for new theorems. It is a
specification --- chapter by chapter, priority by priority ---
for the expository, comparative, and apparatus work needed to
turn a correct manuscript into a readable one. Each priority
section below specifies not only \emph{what} should be added
but \emph{where}, \emph{how much}, and \emph{in what order},
so that the plan itself can be executed incrementally without
re-deriving scope at each session.

\section*{How to read this document}

Each priority (A through H) is self-contained and can be
executed independently of the others, but Section~\ref{sec:sequencing}
gives a recommended execution order based on dependency and
leverage: some priorities (in particular Priority~A, chapter
introductions) make every other priority easier to execute
afterward, because they establish the narrative voice that
extended examples and literature discussions should match.

Throughout, page estimates assume the manuscript's current
typesetting (11pt, LuaLaTeX, standard book margins, the
existing \texttt{tcolorbox}-based theorem environments) and are
based on the current line-count census in
Table~\ref{tab:chapter-census}.

\chapter{Current State Assessment}
\label{ch:assessment}

\section{Manuscript Census}

Table~\ref{tab:chapter-census} gives the current state of every
chapter, including its part, approximate current length, and a
status flag for each of the eight priorities defined in this
plan. A checkmark indicates the priority is substantially
addressed already; a tilde indicates partial coverage; a cross
indicates the priority has not yet been addressed for that
chapter.

\begin{longtable}{@{}p{0.5cm}p{4.3cm}p{1.3cm}cccccccc@{}}
\toprule
\# & Chapter & Lines & A & B & C & D & E & F & G & H \\
\midrule
\endhead
1  & The Problem of Difference        & 407 & \statuspartial & \statusgap & \statuspartial & --- & --- & --- & --- & \statuspartial \\
2  & Distinction and Information      & 215 & \statusgap & \statusgap & \statuspartial & --- & --- & --- & --- & \statusgap \\
3  & Distinction and Entropy          & 252 & \statusgap & \statusgap & \statuspartial & --- & --- & --- & --- & \statusgap \\
4  & History Before State             & 233 & \statusgap & \statusgap & \statusgap & --- & --- & --- & --- & \statusgap \\
5  & Recoverability                   & 234 & \statuspartial & \statuspartial & \statusgap & --- & --- & --- & --- & \statuspartial \\
6  & Memory and Reconstruction        & 284 & \statusgap & \statusgap & \statusgap & --- & --- & --- & --- & \statusgap \\
7  & Repair as a Fundamental Op.\     & 162 & \statusgap & \statusgap & \statusgap & --- & --- & --- & --- & \statusgap \\
8  & The Geometry of Failure          & 172 & \statusgap & \statusgap & \statuspartial & --- & --- & --- & --- & \statusgap \\
9  & Repair and Intelligence          & 180 & \statusgap & \statusgap & \statuspartial & --- & --- & --- & --- & \statusgap \\
10 & Beyond Continuation              & 323 & \statuspartial & \statuspartial & \statusgap & --- & --- & --- & --- & \statusgap \\
11 & Regenerative Systems             & 358 & \statuspartial & \statuspartial & \statusgap & --- & --- & --- & --- & \statusgap \\
12 & Distinction Ecology              & 369 & \statuspartial & \statuspartial & \statusgap & --- & --- & --- & --- & \statuspartial \\
13 & Reachability                     & 584 & \statuspartial & \statuspartial & \statusgap & --- & --- & --- & --- & \statuspartial \\
14 & Admissibility                    & 553 & \statuspartial & \statuspartial & \statusgap & --- & --- & --- & --- & \statuspartial \\
15 & Geometry of Admissible Futures   & 159 & \statusgap & \statusgap & \statusgap & --- & --- & --- & --- & \statusgap \\
16 & RSVP                             & 583 & \statusgap & \statusgap & \statusgap & \statusgap & --- & --- & --- & \statusgap \\
17 & Gravity as Distinction Dynamics  & 499 & \statusgap & \statusgap & \statusgap & \statusgap & --- & --- & --- & \statusgap \\
18 & Expyrotic Cosmology              & 428 & \statusgap & \statusgap & \statusgap & \statusgap & --- & --- & --- & \statusgap \\
19 & Cosmological Completion          & 394 & \statusgap & \statusgap & \statusgap & \statusgap & --- & --- & --- & \statusgap \\
20 & Semantic Geometry                & 687 & \statusgap & \statuspartial & \statusgap & \statusgap & --- & --- & --- & \statuspartial \\
21 & Consciousness                    & 466 & \statusgap & \statusgap & \statusgap & \statusgap & --- & --- & --- & \statusgap \\
22 & Preference Fields                & 442 & \statusgap & \statusgap & \statusgap & \statusgap & --- & --- & --- & \statusgap \\
23 & Flow Computing                   & 323 & \statusgap & \statuspartial & \statusgap & \statusgap & --- & --- & --- & \statuspartial \\
24 & Spherepop                        & 392 & \statusgap & \statuspartial & \statusgap & \statusgap & --- & --- & --- & \statuspartial \\
25 & HYDRA                            & 276 & \statusgap & \statusgap & \statusgap & \statusgap & --- & --- & --- & \statusgap \\
26 & TARTAN                           & 386 & \statusgap & \statuspartial & \statusgap & \statusgap & --- & --- & --- & \statuspartial \\
27 & Fiscal Reachability              & 366 & \statusgap & \statuspartial & \statusgap & \statusgap & --- & --- & --- & \statuspartial \\
28 & Governance as Future Pres.\      & 337 & \statusgap & \statuspartial & \statusgap & \statusgap & --- & --- & --- & \statuspartial \\
29 & Science as Regenerative System   & 423 & \statusgap & \statuspartial & \statusgap & \statusgap & --- & --- & --- & \statuspartial \\
30 & The Ecology of Distinctions      & 826 & \statuspartial & \statusgap & \statusgap & --- & \statuspartial & --- & --- & \statuspartial \\
31 & The Preservation of Possibility  & 735 & \statuspartial & \statusgap & \statusgap & --- & --- & \statuspartial & --- & \statuspartial \\
A  & Mathematical Prerequisites       & 191 & --- & --- & --- & --- & --- & --- & \statuspartial & --- \\
B  & Computational Tools              & 148 & --- & --- & --- & --- & --- & --- & \statuspartial & --- \\
C  & Biological Background            & 137 & --- & --- & --- & --- & --- & --- & \statuspartial & --- \\
D  & Proof Dependency Tables          & 253 & --- & --- & --- & --- & --- & --- & \statuspartial & --- \\
\bottomrule
\caption{Chapter-by-chapter status against the eight editorial
  priorities. A=Introductions, B=Examples, C=Literature,
  D=Realization framing (Parts~V--IX only), E=Synthesis
  expansion (Ch.~30 only), F=Final chapter expansion (Ch.~31
  only), G=Apparatus (appendices only), H=Indexing.}
\label{tab:chapter-census}
\end{longtable}

\section{Reading the Census}
\label{sec:assessment-patterns}

Three patterns are visible in Table~\ref{tab:chapter-census}.

\textbf{First}, Priority~A (chapter introductions) is at best
partially addressed anywhere in the manuscript, and entirely
absent in Parts~III through~IX (Chapters~7--29). Every chapter
opens with an epigraph and, where added in recent sessions, an
objectives box, but no chapter currently contains the
three-to-six pages of motivating prose specified in
Priority~A below.

\textbf{Second}, the realization chapters (16--29, spanning
Physical, Cognitive, Computational, and Social Realizations)
are simultaneously the least introduced, the least exemplified,
and the least explicitly connected back to the abstract theory.
This is the leverage point: these fourteen chapters are where a
reader is most likely to either become convinced the framework
is doing real work, or conclude it is relabeling existing
ideas. Priorities~A, B, and D are jointly most urgent here.

\textbf{Third}, indexing (Priority~H) is partial in exactly the
chapters that already received attention in recent sessions
(5, 12--14, 20, 23--24, 26--29) and absent everywhere else,
confirming that index entries have so far been added
opportunistically rather than systematically. Priority~H
requires a single dedicated pass over the whole document rather
than incremental chapter-by-chapter work.

\chapter{Priority A --- Chapter Introductions}
\label{ch:priority-a}

\begin{prioritybox}{priorityA}{Priority A: Scope}
\textbf{Target:} Every one of the 31 numbered chapters.\\
\textbf{Length:} 3--6 pages of conceptual exposition per
  chapter, inserted between the epigraph/objectives block and
  the first formal definition.\\
\textbf{Estimated total addition:} 31 chapters $\times$
  4.5 pages average $\approx$ 140 pages.\\
\textbf{Most urgent subset:} Chapters 16--29 (zero current
  coverage).
\end{prioritybox}

\section{The Four-Question Template}

Every chapter introduction should answer, in prose, the
following four questions, in this order. The template is
deliberately uniform across all 31 chapters so that a reader
who has internalized it in Chapter~1 can navigate Chapter~25
the same way.

\begin{enumerate}[label=\textbf{Q\arabic*.}, leftmargin=2.2em]
\item \textbf{What problem is being solved?} State the
  specific question the chapter answers, in non-technical
  language, before any notation appears. For Chapter~16
  (RSVP), this is: \textit{what is the minimal physical field
  content needed for a universe whose fundamental quantity is
  distinguishability rather than position or momentum?}
\item \textbf{Why have previous approaches been insufficient?}
  Identify the specific limitation of the field's standard
  treatment that motivates a reformulation. This is where
  Priority~C (literature integration) and Priority~A overlap:
  the introduction names the comparison; the literature section
  develops it.
\item \textbf{Why does the proposed framework emerge naturally
  from the previous chapter?} This is the continuity
  requirement: every chapter introduction should contain at
  least one explicit backward-pointing sentence (e.g.,
  ``Chapter~15 established that admissible futures form a
  manifold $\adm(x,t)$; this chapter asks what physical field
  realizes that manifold'') so that the book reads as a single
  argument rather than 31 independent essays.
\item \textbf{What should the reader expect to learn?} A short,
  concrete preview --- not the objectives box restated, but a
  narrative compression of it (``by the end of this chapter,
  gravity will appear not as a fundamental force but as the
  gradient of a recoverability field'').
\end{enumerate}

\section{Per-Chapter Specification}
\label{sec:priority-a-perchapter}

Table~\ref{tab:priority-a-detail} gives, for each chapter, the
specific Q2 comparison target (the ``previous approach'' the
introduction should position against) and the specific Q3
backward link (the prior chapter or theorem the new material
should be shown to emerge from). These are not optional
flourishes: Q2 and Q3 are the two questions most likely to be
skipped under time pressure, and they are the two that do the
most work in making the manuscript feel like a single argument
rather than a sequence of topics.

\begin{longtable}{@{}p{0.4cm}p{3.6cm}p{5.6cm}p{5.6cm}@{}}
\toprule
\# & Chapter & Q2: Comparison Target & Q3: Backward Link \\
\midrule
\endhead
1 & Problem of Difference & Object-first metaphysics (substance ontology, Aristotelian categories) & --- (opening chapter; link forward instead, to Ch.~2) \\
2 & Distinction \& Information & Shannon information as the primitive (vs.\ distinction as prior to it) & Axiom of Distinction (Ch.~1) \\
3 & Distinction \& Entropy & Boltzmann/Gibbs entropy as physically primitive & First Distinction Theorem (Ch.~1--2) \\
4 & History Before State & State-space dynamics (physics, CS automata theory) & Distinction--Entropy Duality (Ch.~3) \\
5 & Recoverability & Classical information recovery / error correction as purely syntactic & Law of Historical Compression (Ch.~4) \\
6 & Memory \& Reconstruction & Memory as storage (vs.\ memory as preserved recoverability) & Law of Recoverability (Ch.~5) \\
7 & Repair as Fundamental Op. & Optimization as the universal corrective process & Reconstruction Theorem (Ch.~6) \\
8 & Geometry of Failure & Error/anomaly as a scalar quantity (vs.\ a geometric object) & Repair Existence Theorem (Ch.~7) \\
9 & Repair \& Intelligence & Capability-centric definitions of intelligence (benchmarks, task performance) & Ontology Revision Theorem (Ch.~8) \\
10 & Beyond Continuation & Continuation/survival as the success criterion for systems & Intelligence--Repair Equivalence (Ch.~9) \\
11 & Regenerative Systems & Homeostasis as the biological success criterion & Generativity vs.\ continuation distinction (Ch.~10) \\
12 & Distinction Ecology & Resource-ecology models (energy/nutrient flow) & Principle of Regeneration (Ch.~11) \\
13 & Reachability & Reachability in control theory (state-space, not distinction-space) & Distinction Ecology results (Ch.~12) \\
14 & Admissibility & Constraint satisfaction / feasibility (vs.\ future-preserving feasibility) & Future Volume Theorem (Ch.~13) \\
15 & Geometry of Admissible Futures & Phase-space geometry without a future-preservation metric & Admissibility manifold $\adm(x,t)$ (Ch.~14) \\
16 & RSVP & Standard Model field content / GR's stress-energy tensor & Reachability geometry of Part~V \\
17 & Gravity as Distinction Dyn. & General relativity's geometric-mass postulate & RSVP capacity field (Ch.~16) \\
18 & Expyrotic Cosmology & Inflationary cosmology's horizon problem solution & Gravity Realization Theorem (Ch.~17) \\
19 & Cosmological Completion & Heat-death / maximum-entropy endpoints & Reintegration operator (Ch.~18) \\
20 & Semantic Geometry & Distributional/vector-space semantics (word embeddings) & Admissibility geometry (Part~V) \\
21 & Consciousness & Integrated Information Theory, Global Workspace Theory & Recursive self-model machinery \\
22 & Preference Fields & Utility theory / revealed preference & Consciousness's self-model (Ch.~21) \\
23 & Flow Computing & Von Neumann architecture (state-based computation) & Preference flow fields (Ch.~22) \\
24 & Spherepop & Lambda calculus / standard operational semantics & Flow Computing primitives (Ch.~23) \\
25 & HYDRA & Modular cognitive architectures (vs.\ repair-coordinated ones) & Spherepop event calculus (Ch.~24) \\
26 & TARTAN & Symbolic vs.\ connectionist architecture debates & HYDRA's repair coordination (Ch.~25) \\
27 & Fiscal Reachability & Standard public-finance accounting frameworks & TARTAN's architecture (Ch.~26) \\
28 & Governance \& Future Pres. & Social contract theory, mechanism design & Fiscal reachability volume (Ch.~27) \\
29 & Science as Regen.\ System & Popperian falsificationism, Kuhnian paradigms (revisited) & Governance capacity (Ch.~28) \\
30 & Ecology of Distinctions & --- (synthesis chapter; survey all prior comparisons) & All of Parts~I--IX \\
31 & Preservation of Possibility & Existential risk / longtermist frameworks & Ecology of Distinctions hierarchy (Ch.~30) \\
\bottomrule
\caption{Per-chapter comparison targets and backward links for
  Priority~A introductions.}
\label{tab:priority-a-detail}
\end{longtable}

\section{Special Handling for Chapters 16--29}

The realization chapters require a fifth element beyond the
four-question template: an explicit statement, in the
introduction itself, of \emph{what would falsify the claim that
this domain is a genuine instance of the abstract theory} ---
not merely an illustration of it. This anticipates Priority~D
and should be written in coordination with it (see
Chapter~\ref{ch:priority-d}), but the falsifiability statement
belongs in the introduction, before the formal machinery,
because it is what earns the reader's trust that the chapter is
not merely relabeling.

\chapter{Priority B --- Extended Examples}
\label{ch:priority-b}

\begin{prioritybox}{priorityB}{Priority B: Scope}
\textbf{Target:} Every theorem in every chapter; minimum one
  worked example per major theorem (theorems with a dedicated
  \texttt{tcbtheorem} environment and nontrivial proof).\\
\textbf{Placement rule:} Example before proof, wherever the
  example does not depend on results established later in the
  same proof.\\
\textbf{Estimated total addition:} approximately 90--110 worked
  examples across the manuscript, averaging half a page each
  $\approx$ 50 pages.
\end{prioritybox}

\section{The Ten Recommended Examples and Their Homes}

The ten examples named in the original plan already exist in
separate essays in the broader research corpus and can be
integrated directly rather than written from scratch. Each is
assigned below to its most natural chapter, with a note on
which specific theorem it should illustrate.

\begin{longtable}{@{}p{4.3cm}p{2.2cm}p{8.4cm}@{}}
\toprule
Example & Chapter & Illustrates \\
\midrule
\endhead
Cell membranes as maintained distinctions & Ch.~1 & Axiom of Distinction; the cost of \emph{not} drawing a boundary \\
Hash functions as distinction collapse & Ch.~2 & First Distinction Theorem; information loss under projection \\
DNA repair as recoverability restoration & Ch.~5--7 & Law of Recoverability (Ch.~5); Repair Existence Theorem (Ch.~7) \\
Scientific revolutions as ontology repair & Ch.~8, Ch.~29 & Ontology Revision Theorem (Ch.~8); revisited as a regenerative-system case study in Ch.~29 \\
Neural criticality as regenerative balancing & Ch.~11, Ch.~21 & Principle of Regeneration (Ch.~11); connect to Consciousness chapter's recoverability floor \\
Economic institutions as future-preservation systems & Ch.~14, Ch.~28 & Admissibility (Ch.~14); Governance as Future Preservation (Ch.~28) \\
File systems as histories rather than states & Ch.~4 & History Before State; concrete CS-native example for the chapter's central claim \\
Biological evolution as admissibility exploration & Ch.~14--15 & Admissibility manifold; future volume as fitness landscape reinterpretation \\
Bureaucratic collapse as reachability contraction & Ch.~13, Ch.~28 & Reachability (Ch.~13); governance collapse case study \\
Scientific communities as distinction ecologies & Ch.~12, Ch.~29 & Distinction Ecology (Ch.~12); Science as Regenerative System (Ch.~29) \\
\bottomrule
\caption{Mapping of the ten recommended examples to chapters
  and theorems.}
\label{tab:priority-b-examples}
\end{longtable}

\section{Coverage Beyond the Ten Named Examples}

The ten examples above cover roughly a third of the
manuscript's major theorems. The remaining chapters --- in
particular the Physical Realizations (16--19), where the
``example'' is typically a worked numerical or limiting case
rather than a cross-domain analogy --- need a different style
of example:

\begin{itemize}
\item \textbf{Chapters 16--19 (RSVP, Gravity, Expyrosis,
  Cosmology):} worked limiting cases. For example, Chapter~16's
  Energy Decay Theorem should be illustrated with an explicit
  one-dimensional $(\Phi, v, S)$ configuration where $E[t]$ can
  be computed in closed form, so the reader sees the Lyapunov
  argument operate on numbers before trusting it in the
  general case.
\item \textbf{Chapters 20--22 (Semantic Geometry, Consciousness,
  Preference Fields):} small finite-state toy systems. A
  five-state semantic space is sufficient to illustrate the
  Semantic Geometry chapter's central distance theorems without
  requiring the reader to track an infinite-dimensional space.
\item \textbf{Chapters 23--26 (Computational Realizations):}
  these chapters already contain implicit examples in the form
  of pseudocode and operational-semantics derivations (Spherepop
  in particular); Priority~B's job here is to add a short
  worked trace --- a concrete sequence of Pop/Refuse/Collapse/
  Bind events --- before the formal semantics, not to invent new
  examples from scratch.
\item \textbf{Chapters 27--29 (Social Realizations):} these
  chapters benefit most from real-world historical cases (a
  specific historical fiscal crisis, a specific governance
  collapse, a specific scientific-community case study) rather
  than toy examples, since the chapters' claims are explicitly
  about real institutions.
\end{itemize}

\chapter{Priority C --- Literature Integration}
\label{ch:priority-c}

\begin{prioritybox}{priorityC}{Priority C: Scope}
\textbf{Target:} A dedicated ``Related Frameworks'' subsection
  near the end of each chapter (after the main theorems, before
  the Chapter Summary), discussing 1--3 of the thirteen listed
  literatures per chapter as relevant.\\
\textbf{Stance:} Comparative, not subordinating. Each
  discussion should state both what the cited framework gets
  right that this book inherits, and where this book's
  commitments diverge.\\
\textbf{Estimated total addition:} 31 subsections
  $\times$ 1--2 pages $\approx$ 40--50 pages.
\end{prioritybox}

\section{Assignment of Literatures to Chapters}

The thirteen literatures named in the original plan are not
equally relevant to every chapter. Table~\ref{tab:priority-c}
assigns each literature to its 2--4 most natural chapters, so
that the comparative discussion in each chapter is focused
rather than a generic survey repeated 31 times.

\begin{longtable}{@{}p{4cm}p{10.9cm}@{}}
\toprule
Literature & Primary chapters \\
\midrule
\endhead
Shannon information theory & Ch.~2 (Distinction \& Information), Ch.~3 (Distinction \& Entropy) \\
Spencer-Brown's \textit{Laws of Form} & Ch.~1 (already epigraph-linked; deepen the comparison), Ch.~2 \\
Bateson's difference theory & Ch.~1, Ch.~4 (History Before State) \\
Cybernetics (Wiener, Ashby) & Ch.~7 (Repair), Ch.~11 (Regenerative Systems), Ch.~28 (Governance) \\
Systems theory (Luhmann, von Bertalanffy) & Ch.~12 (Distinction Ecology), Ch.~28 (Governance), Ch.~29 (Science) \\
Process philosophy (Whitehead) & Ch.~4 (History Before State), Ch.~10 (Beyond Continuation) \\
Information geometry & Ch.~13 (Reachability), Ch.~14 (Admissibility), Ch.~20 (Semantic Geometry) \\
Dynamical systems theory & Ch.~13 (Reachability), Ch.~16 (RSVP), Ch.~19 (Cosmological Completion) \\
Active inference (Friston) & Ch.~21 (Consciousness), Ch.~22 (Preference Fields) \\
Complexity science & Ch.~11 (Regenerative Systems), Ch.~12 (Distinction Ecology), Ch.~26 (TARTAN) \\
Ecological resilience theory (Holling) & Ch.~12 (Distinction Ecology), Ch.~29 (Science as Regen.\ System) \\
Computational mechanics (Crutchfield) & Ch.~4 (History Before State), Ch.~6 (Memory \& Reconstruction), Ch.~20 (Semantic Geometry) \\
Category theory & Ch.~30 (Ecology of Distinctions: $\mathbf{Reg}$ category), Ch.~31 (Preservation of Possibility: $\mathbf{Reg}_\Pi$ category) \\
\bottomrule
\caption{Assignment of the thirteen literatures to their
  primary chapters.}
\label{tab:priority-c}
\end{longtable}

\section{Required Structure of Each Comparative Discussion}

To keep ``placement, not dependence'' (the original plan's
phrase) operative rather than aspirational, each Related
Frameworks subsection should follow a fixed three-part
structure:

\begin{enumerate}
\item \textbf{Convergence.} What does the cited framework
  identify correctly that this book's machinery formalizes
  differently? (E.g., Shannon correctly identifies that
  information is about reduction of uncertainty; this book's
  contribution is locating the uncertainty-reducing act
  --- distinction --- as prior to the probability space Shannon
  presupposes.)
\item \textbf{Divergence.} What specific claim does this book
  make that the cited framework does not, or could not, make in
  its own terms? This is the substantive content of the
  subsection and should reference a specific theorem by label.
\item \textbf{Open question.} Where the comparison is genuinely
  unsettled (e.g., the relationship between active inference's
  free-energy minimization and this book's admissibility
  preservation in Ch.~22), say so explicitly rather than
  asserting priority.
\end{enumerate}

\chapter{Priority D --- Realization Chapters: ``Why This Realization Matters''}
\label{ch:priority-d}

\begin{prioritybox}{priorityD}{Priority D: Scope}
\textbf{Target:} Chapters 16--29 (Parts~V through~IX: Physical,
  Cognitive, Computational, and Social Realizations) --- 14
  chapters.\\
\textbf{Placement:} A dedicated section, titled exactly
  \textit{Why This Realization Matters}, placed immediately
  after the chapter's central realization theorem (e.g., the
  Gravity Realization Theorem in Ch.~17) and before the Chapter
  Summary.\\
\textbf{Length:} 1.5--3 pages per chapter $\approx$
  25--35 pages total.
\end{prioritybox}

\section{The Genuine-Instance Test}

Every ``Why This Realization Matters'' section must pass what
this plan calls the \emph{genuine-instance test}: it must
identify a specific, falsifiable consequence that follows
\emph{only} if the realization is a true instance of the
abstract theory and not merely dressed-up terminology. If no
such consequence can be named, that is itself an important
editorial finding and should be flagged rather than papered
over with assertive language.

\section{Specification for the Five Named Chapters}

The original plan names five chapters explicitly. Each is
expanded below into a concrete drafting brief.

\subsection*{RSVP (Chapter 16)}
\textbf{Claim to defend:} capacity, transport, and constraint
are not three independently chosen fields but the unique
minimal field content that can realize a reachability geometry
($\Phi e^{-S}$ as a volume measure, $\mathbf{v}$ as its
transport).\\
\textbf{Genuine-instance test:} the Local Well-Posedness
Theorem (\texttt{thm:rsvp-well-posed} in the existing text)
should be shown to fail, or to require additional structure, if
any one of the three fields is removed --- this is the
falsifiable content that distinguishes ``RSVP instantiates
reachability geometry'' from ``RSVP uses similar vocabulary to
reachability geometry.''

\subsection*{Consciousness (Chapter 21)}
\textbf{Claim to defend:} recursive self-reconstruction is not
analogous to repair but is a special case of it, in the formal
sense of satisfying the Repair Operator definition
(\texttt{defn:repair-op}) with $\mathcal{D}$ restricted to the
self-model.\\
\textbf{Genuine-instance test:} the Reconstruction Persistence
Theorem (\texttt{thm:reconstruction-persistence}) should be shown
to be a direct corollary of the general Repair Existence
Theorem (\texttt{thm:repair-exist}) once the restriction to
self-models is made explicit, with the derivation shown rather
than asserted.

\subsection*{Preference Fields (Chapter 22)}
\textbf{Claim to defend:} optimization objectives are
insufficient without a future-preservation constraint, because
an unconstrained optimizer can satisfy any utility function
while monotonically destroying admissible future volume.\\
\textbf{Genuine-instance test:} construct (or reference, if
already present via the Boundary-Flux Theorem,
\texttt{thm:boundary-flux}) an explicit utility function that is
maximized along a trajectory with strictly decreasing
$\Vol(\adm_t)$, demonstrating that utility maximization alone
does not entail the Preference--Admissibility Equivalence
(\texttt{thm:pref-adm-equivalence}).

\subsection*{HYDRA (Chapter 25)}
\textbf{Claim to defend:} intelligence, as formalized in
Chapter~9's Intelligence--Repair Equivalence
(\texttt{thm:intel-repair}), emerges in HYDRA specifically from
\emph{coordination} between repair processes, not from any
single component's capacity.\\
\textbf{Genuine-instance test:} identify a HYDRA configuration
in which every individual module has positive repair capacity
$\kappa$ but the coordinated system's $\mathcal{I}(\Sigma)$ is
lower than the sum of components --- demonstrating that
coordination failure, not component failure, is what the
realization predicts and that this prediction is specific to
the repair-coordination reading of intelligence.

\subsection*{Fiscal Reachability (Chapter 27)}
\textbf{Claim to defend:} public finance is fundamentally
geometric (a question of which future fiscal states remain
reachable) rather than accounting-based (a question of current
balance).\\
\textbf{Genuine-instance test:} construct a scenario with
identical current-period accounting balance but different
$\Vol(\adm_t)$ under the chapter's fiscal reachability
machinery, and show the two scenarios are treated identically
by standard accounting frameworks but differently --- and, the
section should argue, correctly differently --- by this
chapter's geometric treatment.

\section{Specification for the Remaining Nine Chapters}

Chapters 17--20, 23--24, 26, and 28--29 require the same
treatment; Table~\ref{tab:priority-d-remaining} gives the claim
to defend and a candidate genuine-instance test for each,
sufficient to begin drafting without further scoping.

\begin{longtable}{@{}p{3.6cm}p{6cm}p{5.3cm}@{}}
\toprule
Chapter & Claim to Defend & Candidate Test \\
\midrule
\endhead
17: Gravity & Gravity is the geometric manifestation of capacity gradients, not an independent force law & Show the Distinction Raychaudhuri equation (\texttt{thm:dist-raychaudhuri}) reduces to the standard Raychaudhuri equation under a stated limit, with the limit's physical meaning made explicit \\
18: Expyrotic Cosmology & Cosmic renewal is a repair operation on the global distinction manifold, not a separate cosmological postulate & Show the Reintegration Fixed-Point Theorem (\texttt{thm:reintegration-fixed-point}) is structurally identical to the Repair Existence Theorem (\texttt{thm:repair-exist}) under the appropriate domain restriction \\
19: Cosmological Completion & ``Heat death'' and ``saturation without exhaustion'' are distinguishable, falsifiably different end states & Contrast the Asymptotic Saturation Theorem (\texttt{thm:asymptotic-saturation}) against a heat-death model on the same field content, identifying the divergent prediction \\
20: Semantic Geometry & Meaning is a geometric property of the admissibility manifold, not a property of token co-occurrence statistics & Identify a semantic distinction the chapter's geometry predicts but distributional semantics cannot represent without additional machinery \\
23: Flow Computing & Computation is fundamentally a flow process on distinctions, not state transformation & Identify a computational property invisible to state-based semantics but visible under the flow formalization \\
24: Spherepop & The Pop/Refuse/Collapse/Bind event calculus is the natural operational semantics for irreversible distinction-events, not an arbitrary alternative to lambda calculus & Show a program whose irreversibility is naturally expressed in Spherepop but requires extra-linguistic annotation in standard operational semantics \\
26: TARTAN & The architecture's symbolic/connectionist integration is a repair-coordination structure, not an engineering compromise & Apply the same coordination-failure test used for HYDRA (Ch.~25) to TARTAN's specific architecture \\
28: Governance & Governance is fundamentally about preserving future reachable states for a polity, not about present-period institutional design & Construct an institution that is well-designed by present-period criteria but reduces $\Vol(\adm_t)$ for the polity, and one that is the converse \\
29: Science & Scientific practice is a regenerative system in the chapter's formal sense, not merely analogous to one & Show that the Scientific Revolution Theorem (\texttt{thm:sci-rev-thm}, Ch.~8) and the chapter's own regenerative-system machinery make the same prediction about paradigm-shift timing, derived independently \\
\bottomrule
\caption{Genuine-instance tests for the remaining nine
  realization chapters.}
\label{tab:priority-d-remaining}
\end{longtable}

\chapter{Priority E --- Synthesis Expansion (Chapter 30)}
\label{ch:priority-e}

\begin{prioritybox}{priorityE}{Priority E: Scope}
\textbf{Target:} Chapter~30, ``The Ecology of Distinctions''
  (currently 826 lines, the manuscript's second-longest
  chapter).\\
\textbf{Reframing:} from summary chapter to unification
  chapter. The distinction matters: a summary chapter restates
  results already proved; a unification chapter proves new
  results \emph{about} the relationships between prior results.\\
\textbf{Estimated addition:} 30--45 pages, given the seven
  content areas below.
\end{prioritybox}

\section{Current Coverage}

Chapter~30 already contains the Preservation Hierarchy Theorem
(\texttt{thm:preservation-hierarchy}) and, as of the most recent
revision, the Structural Dependency Theorem
(\texttt{thm:structural-dependency}) establishing the collapse
direction of that hierarchy. These are unification results in
the intended sense and should be the model for the seven
additions below, not superseded by them.

\section{Seven Content Areas}

\begin{enumerate}
\item \textbf{Distinction conservation.} A formal conservation
  law for $D$ analogous to the Distinction Balance Law already
  present in Chapter~31 (\texttt{thm:dist-balance}), but derived
  here from first principles using only Parts~I--II machinery,
  establishing that the balance law is not assumed but follows
  from the Axiom of Distinction.
\item \textbf{Repair conservation.} The analogous conservation
  law for repair capacity $\kappa_\repair$, connecting the
  Repair Conservation Law of Chapter~7
  (\texttt{thm:repair-conservation}, currently a statement about
  individual repairs preserving connected components of
  $\mathcal{M}_\reco$) to an aggregate, ecology-wide version.
\item \textbf{Reachability conservation.} A conservation law
  for $V_F$ under composition of admissible morphisms, proved
  using the Future Volume Theorem (Ch.~13) together with the
  Boundary-Flux Theorem (Ch.~22), making explicit that
  reachability is conserved under the same conditions that make
  preference fields generatively admissible.
\item \textbf{Admissibility conservation.} The corresponding
  statement for $\Vol(\adm)$ itself, which should be shown to
  be the special case of reachability conservation obtained by
  restricting to admissible trajectories --- making explicit
  the relationship between Chapters~13 and~14 that is currently
  left implicit.
\item \textbf{Functorial relationships between realization
  domains.} A formal treatment of the maps between the
  Physical, Cognitive, Computational, and Social Realizations
  (Parts~V--IX), framed as functors between categories of
  realizations that preserve the abstract admissibility
  structure. This is the single most significant addition in
  Priority~E: it would make precise, for the first time in the
  manuscript, the sense in which RSVP (Ch.~16) and Fiscal
  Reachability (Ch.~27) are ``the same theory in different
  domains'' rather than merely thematically similar.
\item \textbf{The category of regenerative systems.} Chapter~30
  already defines $\mathbf{Reg}$ (\texttt{thm:regen-category});
  this addition should explore its categorical properties in
  more depth --- limits, colimits, and whether $\mathbf{Reg}$
  has a terminal or initial object corresponding to maximal or
  minimal regenerative capacity.
\item \textbf{Universality classes of distinction ecologies.}
  A classification result, in the spirit of universality
  classes in dynamical systems theory (cf.\ Priority~C's
  literature assignment), grouping distinction ecologies by
  their long-run qualitative behavior (e.g., bounded
  regeneration, unbounded growth, collapse) rather than by
  domain.
\end{enumerate}

\section{Sequencing Within Priority E}

Items 1--4 (the four conservation laws) should be drafted
first, since item 5 (functorial relationships) depends on
having all four conserved quantities available as the data a
functor must preserve. Item 6 builds directly on item 5,
since functoriality is most naturally stated in terms of
morphisms in $\mathbf{Reg}$. Item 7 should be drafted last, as
a genuine open research question rather than a fully resolved
result --- consistent with the spirit of Priority~F's
recommendation that Chapter~31 read more like a research
program than a closed theorem.

\chapter{Priority F --- Final Chapter Expansion (Chapter 31)}
\label{ch:priority-f}

\begin{prioritybox}{priorityF}{Priority F: Scope}
\textbf{Target:} Chapter~31, ``The Preservation of
  Possibility'' (currently 735 lines, the manuscript's longest
  chapter after Chapter~30).\\
\textbf{Reframing:} from a chapter that states the Generative
  Admissibility Principle to one that mounts a sustained
  normative argument for it.\\
\textbf{Estimated addition:} 20--30 pages.
\end{prioritybox}

\section{Current Coverage}

Chapter~31 already contains the formal apparatus: the refined
possibility functional $\Pi$, the Preservation Implication and
Preservation Objective Equivalence Theorems, the
$\mathbf{Reg}_\Pi$ category, the Unified Invariant Theorem, and
the Generative Domination Theorem. What it does not yet contain
is a sustained discussion of \emph{why} a reader should accept
possibility-preservation as the relevant normative criterion,
as opposed to treating it as one formally elegant option among
several. The seven topics below are precisely this missing
normative layer.

\section{Seven Topics, Each as a Comparative Argument}

Each topic should be structured as a direct comparison: name
the rival criterion, state its strongest form, and show
specifically where it diverges from possibility-preservation
and why that divergence favors the latter — or, where it does
not clearly favor either, say so.

\begin{enumerate}
\item \textbf{Why continuation is insufficient.} The chapter
  already has the materials for this: the Survival--Viability
  Separation Theorem (\texttt{thm:survival-viability}) shows
  systems can survive indefinitely while $\mathcal{P}\to0$.
  This topic should expand the \emph{interpretation} of that
  theorem --- what does it mean, normatively, for a system to
  satisfy the weakest possible success criterion while losing
  everything of value by this book's lights?
\item \textbf{Why survival is insufficient.} Distinct from
  continuation: a system can fail to continue (in the sense of
  Ch.~10) while nonetheless transferring possibility elsewhere
  (e.g., an organism's death transferring distinction-capacity
  to offspring or ecosystem). This topic should use the
  Generative Domination Theorem
  (\texttt{thm:generative-domination}) to argue that
  possibility-dominance, not individual survival, is the
  coherent unit of normative evaluation.
\item \textbf{Why optimization is insufficient.} This connects
  directly to Priority~D's Preference Fields treatment
  (Ch.~22): an optimizer can maximize any fixed objective while
  the Boundary-Flux Theorem shows admissible volume can
  simultaneously decrease. This topic should generalize that
  specific result into the chapter's central normative claim.
\item \textbf{Why future-preservation appears repeatedly across
  domains.} This is the explanatory burden the
  Preservation Objective Equivalence Theorem
  (\texttt{thm:preservation-objective-equivalence}) carries
  formally; this topic should carry it normatively, arguing
  that the recurrence is not coincidental convergence on a
  useful heuristic but evidence that possibility-preservation
  is the structurally correct criterion of which the
  domain-specific criteria (repair capacity, institutional
  resilience, etc.) are projections.
\item \textbf{How admissibility differs from utility
  maximization.} A direct technical and philosophical
  comparison: utility maximization is agent-relative and
  requires a specified utility function; admissibility
  preservation is defined without reference to any
  particular agent's preferences, only to the structure of
  future possibility itself. This topic should state precisely
  what is gained and lost by this generality.
\item \textbf{Limitations of the framework.} This topic is
  required, not optional, for intellectual honesty. Candidate
  limitations to address: (i) the framework's reliance on a
  well-defined measure $\mu_R$ for computing $\Vol(\adm)$,
  which may not exist or be unique in all domains; (ii) the
  computational intractability of evaluating $\Pi$ or
  $\mathcal{U}$ in realistic systems; (iii) the open question
  of how to handle genuine moral disagreement about which
  futures count as ``admissible'' in social domains (Ch.~28).
\item \textbf{Open problems.} A closing list, each problem
  stated precisely enough to be picked up independently:
  candidates include extending the Unified Invariant Theorem to
  systems with negative or contested $S_A$; characterizing the
  topology of $\mathbf{Reg}_\Pi$; and determining whether the
  Generative Admissibility Principle admits a
  game-theoretic formulation for multi-agent ecologies.
\end{enumerate}

\section{Tone Requirement}

The original plan specifies that this chapter ``should feel
less like a theorem and more like a research program.'' This
has a concrete drafting consequence: topics 1--5 should be
written with full theorem-citation rigor (every normative claim
anchored to a specific result already proved), while topics 6
and 7 should be written in a more exploratory register, with
explicit markers (``it is not yet known whether\ldots'',
``one natural conjecture is\ldots'') that are largely absent
elsewhere in the manuscript but appropriate here.

\chapter{Priority G --- Scholarly Apparatus (Appendices)}
\label{ch:priority-g}

\begin{prioritybox}{priorityG}{Priority G: Scope}
\textbf{Target:} All four appendices (currently 191, 148, 137,
  and 253 lines respectively --- by far the shortest material
  in the manuscript relative to its apparatus role).\\
\textbf{Estimated addition:} 60--90 pages across the four
  appendices.
\end{prioritybox}

\section{Appendix A: Mathematical Prerequisites \texorpdfstring{$\to$}{->} Compact Handbook}

Currently 191 lines. Target structure as a genuine handbook:

\begin{itemize}
\item Measure-theoretic preliminaries used throughout (the
  measure $\mu_R$ underlying $\Vol(\adm)$, deserving a
  self-contained treatment given how load-bearing it is across
  Chapters~13--31).
\item A consolidated notation table (currently notation is
  introduced locally per chapter; a single table mapping every
  recurring symbol --- $\Phi$, $S$, $\reco$, $\adm$, $\Pi$,
  $\mathcal{U}$, \ldots\ --- to its defining chapter and
  equation would materially reduce the burden on a reader
  jumping between chapters).
\item A short primer on the \texttt{tcbtheorem}-style proof
  conventions used in the book (what counts as a proof sketch
  vs.\ a full proof, and why certain results, e.g.\ the Minimal
  Repair Theorem in Ch.~7, are explicitly marked
  \textit{(Sketch)}).
\end{itemize}

\section{Appendix B: Computational Tools \texorpdfstring{$\to$}{->} Implementation Notes and Formal Semantics}

Currently 148 lines. This appendix should absorb and formalize
the operational content already implicit in the Computational
Realizations (Part~VII, Chapters~23--26):

\begin{itemize}
\item Formal operational semantics for Spherepop's
  Pop/Refuse/Collapse/Bind calculus (Ch.~24), stated in
  standard small-step style, cross-referenced to but not
  duplicating the chapter's own treatment.
\item Implementation notes connecting the book's formalism to
  the existing \texttt{cprsh} runtime and \texttt{spherepop}
  interpreter referenced in the broader research corpus, with
  an explicit statement of which formal results in Chapters
  23--26 are mechanically verified by the existing test suites
  and which remain to be checked.
\item A worked translation of one complete Flow Computing
  (Ch.~23) program into Spherepop (Ch.~24) and into TARTAN
  (Ch.~26) representations, demonstrating the equivalence
  claims made informally across those three chapters.
\end{itemize}

\section{Appendix C: Biological Background \texorpdfstring{$\to$}{->} Biological Realizations of Repair and Regeneration}

Currently 137 lines, the shortest appendix. Target additions:

\begin{itemize}
\item DNA repair mechanisms (nucleotide excision repair,
  mismatch repair) as concrete instances of the Repair Operator
  definition (Ch.~7), extending the worked example already
  assigned to that chapter under Priority~B.
\item Immune memory as a biological instance of the Memory and
  Reconstruction machinery (Ch.~6), with explicit mapping of
  immunological terms (antigen, memory cell) onto the chapter's
  formal vocabulary ($d$, $\reco(d)$, reconstruction operator).
\item Neural criticality and homeostatic plasticity as the
  biological grounding for the Regenerative Systems chapter's
  (Ch.~11) abstract machinery, extending the worked example
  already assigned there under Priority~B.
\item Regeneration in planaria and other classical regeneration
  model organisms, as the most literal biological instance of
  the Principle of Regeneration (Ch.~11) --- a case where
  ``regeneration'' is not a metaphor extension but the original
  meaning of the term the book's vocabulary borrows.
\end{itemize}

\section{Appendix D: Proof Dependency Tables \texorpdfstring{$\to$}{->} Complete Dependency Atlas}

Currently 253 lines, already the most developed appendix and
the natural target for the most substantial Priority~G
expansion, since it is infrastructure that makes the rest of
the manuscript easier to verify and navigate.

\begin{itemize}
\item Extend the existing dependency tables to cover every
  theorem added in the most recent revision sessions (the
  well-posedness/attractor machinery in Ch.~16, the
  curvature-recoverability correspondence in Ch.~17, the
  Structural Dependency Theorem in Ch.~30, and the full
  apparatus of the revised Ch.~31), none of which are yet
  reflected in the appendix.
\item Add a visual dependency graph (a large foldout or
  multi-page TikZ diagram) showing the complete theorem
  dependency structure of the manuscript, with Parts color-coded,
  complementing the existing textual tables.
\item Add a ``minimal path'' table: for each realization chapter
  (16--29), the shortest chain of theorems from the Axiom of
  Distinction (Ch.~1) to that chapter's central realization
  theorem, so a reader who wants to verify one realization
  chapter's foundations without reading the entire book linearly
  can do so.
\end{itemize}

\chapter{Priority H --- Indexing}
\label{ch:priority-h}

\begin{prioritybox}{priorityH}{Priority H: Scope}
\textbf{Target:} Whole-document pass, all 35 chapters.\\
\textbf{Method:} Single dedicated pass rather than incremental
  per-chapter work (see Section~\ref{sec:assessment-patterns}
  above on why opportunistic indexing has produced uneven
  coverage).\\
\textbf{Mechanism:} \texttt{\textbackslash index\{\}} markers
  at first substantive use of each term per chapter (not every
  occurrence), feeding the existing \texttt{makeidx}/
  \texttt{imakeidx} pipeline already configured via
  \texttt{\textbackslash printindex} in the backmatter.
\end{prioritybox}

\section{Core Index Terms and Required Sub-Entries}

The fourteen terms named in the original plan are the
top-level entries; each requires sub-entries distinguishing its
uses across Parts, since several of these terms (notably
\textit{admissibility}, \textit{distinction}, and
\textit{reachability}) appear in meaningfully different senses
across the abstract chapters (Parts~I--V) and the realization
chapters (Parts~V--IX).

\begin{longtable}{@{}p{3.2cm}p{10.8cm}@{}}
\toprule
Top-level entry & Required sub-entries \\
\midrule
\endhead
admissibility & abstract definition (Ch.~14); physical (Ch.~16--19); cognitive (Ch.~21--22); computational (Ch.~23--26); social (Ch.~27--29); generative vs.\ extractive (Ch.~31) \\
blind spot & definition (Ch.~1); relation to partition refinement (Ch.~1--2) \\
collapse & wavefunction-style usage (none --- flag if accidentally conflated); admissibility collapse (Ch.~19); institutional collapse (Ch.~28) \\
distinction & axiom (Ch.~1); cost (Ch.~1); density (Ch.~17); entropy (Ch.~31) \\
ecology & distinction ecology (Ch.~12); scientific communities as (Ch.~29, via Priority~B); ecology of distinctions (Ch.~30) \\
entropy & distinction-theoretic (Ch.~3); admissible distinction entropy $S_A$ (Ch.~31); RSVP field $S$ (Ch.~16) \\
future cone & definition (Ch.~13); relation to $\adm(x,t)$ (Ch.~14--15) \\
history & vs.\ state (Ch.~4); historical self-information (Ch.~21) \\
ontology revision & definition (Ch.~8); generative admissibility of (Ch.~8); meta-learning connection (Ch.~9) \\
reachability & abstract (Ch.~13); control-theory contrast (Priority~C, Ch.~13); fiscal (Ch.~27) \\
recoverability & law of (Ch.~5); floor / threshold (Ch.~21) \\
regeneration & principle of (Ch.~11); biological (Appendix~C, via Priority~G); cosmological (Ch.~18--19) \\
repair & operator (Ch.~7); conservation (Ch.~7, Ch.~30 via Priority~E); intelligence as (Ch.~9) \\
trajectory & admissible (Ch.~14--15); generative/extractive/neutral classes (Ch.~31) \\
\bottomrule
\caption{Required sub-entry structure for the fourteen core
  index terms.}
\label{tab:priority-h}
\end{longtable}

\section{Additional Terms Warranting Index Entries}

Beyond the fourteen named terms, the following recur often
enough across multiple Parts to warrant indexing, identified by
cross-referencing the chapter census in
Table~\ref{tab:chapter-census}: \textit{capacity field},
\textit{generativity}, \textit{Lyapunov function} (appears in
Ch.~16's energy decay and is referenced informally elsewhere),
\textit{possibility functional}, \textit{preservation
hierarchy}, and \textit{semantic horizon}.

\chapter{Recommended Execution Order}
\label{ch:sequencing}
\label{sec:sequencing}

\section{Dependency-Based Sequencing}

Not all eight priorities are independent in practice, even
though each is independently specified above. The following
order is recommended on the grounds that earlier items reduce
the cost of later ones:

\begin{enumerate}
\item \textbf{Priority A (Introductions) first, across all 31
  chapters.} This establishes the narrative voice and the
  explicit backward-link sentences that Priorities B and D will
  need to reference. Drafting B or D before A risks producing
  examples and realization-arguments that have to be retrofitted
  once the introductions establish how each chapter frames its
  own problem.
\item \textbf{Priority H (Indexing) second, as a single pass.}
  Indexing is cheapest to do once, after a chapter's prose is
  stable, and most expensive to redo --- doing it before
  Priorities B--G would mean re-indexing every chapter after
  each subsequent expansion. Doing it immediately after Priority~A
  (rather than last) means the remaining priorities can be
  indexed incrementally as they are written, rather than requiring
  a second whole-document pass at the very end.
\item \textbf{Priority D (realization framing) third, for
  Chapters 16--29 specifically.} This should follow Priority~A
  (which already establishes the comparison targets these
  chapters need) and precede Priority~B for the same chapters,
  since the ``genuine-instance test'' specified in Priority~D
  often \emph{is} the worked example Priority~B would otherwise
  need to invent separately.
\item \textbf{Priority B (Extended Examples) fourth.} By this
  point the comparison targets (Priority~A) and, for realization
  chapters, the genuine-instance tests (Priority~D) are
  available to inform example selection.
\item \textbf{Priority C (Literature Integration) fifth.} This
  depends on Priority~A's Q2 comparison targets
  (Table~\ref{tab:priority-a-detail}) being finalized, since the
  Related Frameworks subsections are the natural place to
  develop a comparison that the introduction only names.
\item \textbf{Priorities E and F (Synthesis and Final Chapter
  Expansion) sixth, together.} These two chapters should be
  revised in the same pass because Priority~F's normative
  argument (Ch.~31) depends on Priority~E's conservation laws
  and functorial relationships (Ch.~30) being in place first ---
  in particular, Priority~F's topic 4 (``why future-preservation
  appears repeatedly across domains'') is most convincing once
  Priority~E's functorial-relationships work has made the
  cross-domain recurrence formally explicit rather than
  asserted.
\item \textbf{Priority G (Scholarly Apparatus) last.} The
  appendices --- in particular the dependency atlas (Appendix~D)
  and notation table (Appendix~A) --- should be finalized only
  after all preceding priorities are complete, since every one
  of them potentially adds new theorems, new notation, or new
  dependency edges that the apparatus must reflect.
\end{enumerate}

\section{Estimated Total Addition}

\begin{deliverable}
\begin{tabular}{@{}lr@{}}
Priority A (Introductions) & $\approx$140 pages \\
Priority B (Extended Examples) & $\approx$50 pages \\
Priority C (Literature Integration) & $\approx$45 pages \\
Priority D (Realization Framing) & $\approx$30 pages \\
Priority E (Synthesis Expansion) & $\approx$38 pages \\
Priority F (Final Chapter Expansion) & $\approx$25 pages \\
Priority G (Scholarly Apparatus) & $\approx$75 pages \\
Priority H (Indexing) & negligible page count; substantial labor \\
\midrule
\textbf{Total estimated addition} & \textbf{$\approx$400--410 pages} \\
\end{tabular}
\end{deliverable}

Given the manuscript's current length (approximately 13{,}250
lines of source, corresponding to roughly 380--420 typeset
pages at the book's current density), full execution of this
plan would roughly double the manuscript's length. This is
consistent with the plan's own framing: the mathematics is
complete; the prose that makes the mathematics legible is, by
volume, comparable in size to the mathematics itself.

\chapter*{Final Goal}
\addcontentsline{toc}{chapter}{Final Goal}

The remaining challenge is no longer proving the theory. The
remaining challenge is making the theory readable.

A completed manuscript should allow a reader to understand the
central argument even if they skip half of the proofs. The
mathematics establishes correctness. The prose establishes
comprehension. Both are necessary, and this plan exists to make
the second kind of work as tractable, sequenced, and concretely
scoped as the first kind already is.

\end{document}
